\section{引言}
可复用技能将任务说明、脚本、参考资料和约束组织为可供大模型智能体调用的能力单元。
技能网络进一步通过检索、关系扩展与编排，把大量技能连接到任务执行流程。
然而，一个节点存在于目录、文档表述完整、模型愿意给出高分，均不足以证明它在指定条件下可以完成任务。
同样，网络节点增多、连接变密或检索相似度提高，也不能直接推出任务成功率提高。
这使开放技能生态同时面临两个测量问题：怎样核验单个技能的质量及边界，以及怎样依据核验证据改进网络选择。

已有工作已提供多维评分、人工评审、工具执行与任务基准。
SkillNet 将技能创建、评价和连接整合到网络框架中\citep{P01}；
SkillCT 将质量分解为静态、动态和人类辅助维度，并提出通过参考画像估计新技能的免执行基线\citep{P43}；
SkillLearnBench 则区分技能文本、执行轨迹与最终任务结果\citep{P44}。
因此，新增几个评分维度，或者让模型多评一轮，并不能自动构成实质性进展。
关键在于：一个分数究竟由哪类证据支持，这种证据允许作出什么决定，
以及决定是否经独立产物验收证明有用。

本研究在现有 SkillNet-demo 工程基础上开展一组可回放的开发实验。
研究不预设所有技能均可获得统一的永久好坏标签，而把对象固定为“技能包版本、任务条件和使用决定”的组合。
对于可机械判定的声明，使用资源指纹、接口契约与独立提取的来源窗口建立反证；
对于需要行为证据的声明，保持未验证状态，并将实际产物交由与模型分离的验收器判定。
在网络侧，则同时测量单项召回与完整能力集合覆盖，防止平均检索分数掩盖组合任务的缺口。

前期审查发现，文档示例可能与脚本参数声明不一致，部分实现方法也可能直接返回字面成功值。
这些现象说明“声明成功”和“有产物支持的成功”需要分开。
但反向推断也同样危险：一个静态异常不等于整个技能在所有用途上无效；
引用没有按规定完整复制，也不等于模型虚构了来源。
本研究将这些区分写入状态模型，并保留每种测量的分母与未知比例。

围绕上述目标，本文回答三个可检验问题：
\begin{itemize}
\item RQ1：在完整资源登记与已知干预条件下，哪些技能声明可以被机械核验，静态方法的覆盖边界在哪里？
\item RQ2：来源可定位、模型评分与实际产物合格之间能否互相替代？证据约束补充说明对指定任务的实际效用如何？
\item RQ3：在固定候选池和选择数量下，集合多样性是否改善技能网络的必要能力覆盖？
\end{itemize}

本文的贡献是一个分层证据协议及其经验分析，而非对所有技能有效的认证器。
首先，建立声明、静态支持、冲突、执行验收和未知的显式区分，给出可回放的资源与来源绑定。
其次，将真实包审查、受控接口干预、真实模型评价和新构造任务的产物验收连接起来，
用强对照与负结果限定方法的有效范围。
最后，保留全部网络选择对照及能力集合指标，说明在何种意义上系统仍未获得稳定改进。
这些发现为后续独立人审、跨来源验证和真实工具编排提供了可以继续检验的基础。
